\documentclass[11pt,a4paper,oneside]{report}

\usepackage[T1]{fontenc}
\usepackage[utf8]{inputenc}
\usepackage[italian]{babel}
\usepackage[a4paper,margin=2.8cm,top=3cm,bottom=3cm]{geometry}
\usepackage{amsmath}
\usepackage{mathpazo}
\usepackage{xcolor}
\usepackage{titlesec}
\usepackage{booktabs}
\usepackage{longtable}
\usepackage{array}
\usepackage{enumitem}
\usepackage{fancyhdr}
\usepackage{graphicx}
\usepackage[hidelinks]{hyperref}

\definecolor{accent}{HTML}{1F3A5F}
\definecolor{soft}{HTML}{6B7C93}
\definecolor{rule}{HTML}{C8D0DA}
\definecolor{boxbg}{HTML}{F4F6F9}

\hypersetup{colorlinks=true,linkcolor=accent,urlcolor=accent,citecolor=accent}

\titleformat{\chapter}[display]
  {\normalfont\Large\bfseries\color{accent}}
  {\normalfont\normalsize\color{soft}\MakeUppercase{\chaptertitlename}\ \thechapter}
  {0.6em}{\Huge}
\titlespacing*{\chapter}{0pt}{-10pt}{30pt}
\titleformat{\section}{\normalfont\large\bfseries\color{accent}}{\thesection}{0.6em}{}
\titleformat{\subsection}{\normalfont\normalsize\bfseries\color{accent}}{\thesubsection}{0.6em}{}

\pagestyle{fancy}
\fancyhf{}
\renewcommand{\headrulewidth}{0.4pt}
\fancyhead[L]{\small\color{soft}\nouppercase{\leftmark}}
\fancyhead[R]{\small\color{soft}\thepage}

\setlist[itemize]{leftmargin=1.3em,itemsep=2pt,topsep=4pt}
\setlist[enumerate]{leftmargin=1.5em,itemsep=2pt,topsep=4pt}

\newcommand{\metrica}[1]{\textcolor{accent}{\textbf{#1}}}
\newenvironment{nota}
  {\begin{quote}\small\color{soft}\itshape}
  {\end{quote}}

% Verdetto box
\newcommand{\verdetto}[2]{%
  \par\vspace{4pt}%
  \noindent\colorbox{boxbg}{%
    \begin{minipage}{\dimexpr\linewidth-2\fboxsep}%
      \vspace{3pt}\textbf{\color{accent}#1}\par\vspace{2pt}#2\vspace{3pt}%
    \end{minipage}}%
  \par\vspace{6pt}}

\begin{document}

% ------------------------------------------------------------------
\begin{titlepage}
\centering
\vspace*{3cm}
{\color{soft}\large Nota metodologica}\\[1.2cm]
{\color{accent}\Huge\bfseries Valutare modelli ML\\[0.3cm] nel data plane}\\[0.8cm]
{\Large Come lo fa la letteratura,\\ come siamo messi, cosa manca}\\[2.5cm]
\rule{0.6\linewidth}{0.8pt}\\[0.8cm]
{\large Progetto \textbf{IPA/eBPF} --- pipeline hardcoded, template, modular}\\[0.3cm]
{\color{soft} topologia Germany50 su Kathara}\\[3cm]
{\color{soft} Settembre 2026}
\end{titlepage}

\tableofcontents
\clearpage

% ==================================================================
\chapter{Perché questa nota}

Il professore ha posto una domanda precisa: \emph{come vengono effettuati
i test per questi tipi di modelli in letteratura?} La domanda non è
retorica. Il progetto misura già molte cose --- istruzioni eBPF, codice
jited, tail call, letture di mappa, memoria, latenza --- e le misura con
cura. Il punto è capire se l'insieme di quelle misure costituisce una
\emph{valutazione} nel senso in cui la intende la letteratura sul
machine learning nel data plane, oppure se copre solo una parte del
problema.

La risposta breve, anticipata qui e argomentata nel resto del
documento, è che il progetto misura molto bene \textbf{una} delle
quattro dimensioni che la letteratura considera obbligatorie, e non
misura affatto le altre tre. Non è un difetto di esecuzione: le misure
esistenti sono più rigorose della media. È un difetto di
\emph{copertura}.

Questa nota è organizzata così: il Capitolo~\ref{ch:cosa} ricostruisce
quali metriche la letteratura riporta e perché; il
Capitolo~\ref{ch:come} descrive con quali strumenti e protocolli le
ottiene; il Capitolo~\ref{ch:trappole} raccoglie gli errori
metodologici documentati, che sono il motivo per cui certe scelte
vanno giustificate esplicitamente; il Capitolo~\ref{ch:noi} confronta
lo stato del progetto con quel quadro; il Capitolo~\ref{ch:fare}
propone un piano di lavoro ordinato per priorità, con criteri di
accettazione.

\begin{nota}
Tutti i riferimenti citati sono stati verificati: titolo, autori, sede
e anno provengono dalle pagine ufficiali degli editori o dagli archivi
degli autori, non dalla memoria. Dove un dato numerico proviene dal
codice di questo progetto è indicato esplicitamente.
\end{nota}

% ==================================================================
\chapter{Che cosa misura la letteratura}
\label{ch:cosa}

Il campo ha un nome consolidato: \emph{in-network machine learning}.
Due rassegne recenti ne definiscono i confini e ne sistematizzano i
risultati: quella di Zheng et al. su IEEE Communications Surveys \&
Tutorials~\cite{zheng24survey} e quella di Parizotto et al. su ACM
Computing Surveys~\cite{parizotto23survey}. Leggendole insieme ai
lavori sperimentali principali emerge una struttura ricorrente: le
valutazioni si articolano su quattro assi, e i lavori pubblicati li
coprono tutti e quattro.

\section{Asse 1 --- Qualità del modello}

È la prima cosa che ogni lavoro riporta, e in genere occupa la prima
tabella. Non si tratta di prestazioni ma di \emph{correttezza
statistica}: quanto bene il modello svolge il compito per cui è stato
addestrato. Le metriche sono quelle standard del machine learning ---
accuratezza, precision, recall, F1 --- valutate su un dataset di test
tenuto separato dall'addestramento.

Il punto cruciale, e quello più spesso trascurato da chi arriva al
problema dal lato sistemistico, è che questa misura va riportata
\textbf{due volte}: una per il modello originale in virgola mobile e
una per il modello effettivamente eseguito nel dispositivo di rete,
quantizzato e vincolato. La differenza fra le due è un risultato in
sé, perché quantifica quanto si è pagato in qualità per poter eseguire
il modello nel data plane.

N3IC~\cite{siracusano22n3ic} --- già citato nella bibliografia della
tesi --- è esemplare in questo: compila reti neurali \emph{binarie}
nel data plane di SmartNIC, e siccome la binarizzazione è una forma
estrema di quantizzazione, il lavoro è costretto a dimostrare che
l'accuratezza resta accettabile su due casi d'uso reali
(identificazione del traffico e anomaly detection). IIsy~\cite{zheng24iisy}
va oltre e fa della perdita di qualità un elemento di progetto: adotta
un modello ``piccolo'' sullo switch e uno ``grande'' nel backend, e usa
la confidenza della decisione per stabilire quando il modello in rete
basta e quando serve inoltrare a valle.

\verdetto{Che cosa comporta per noi}{Il progetto non misura questo asse
in nessuna forma. L'accuratezza è trattata come \emph{invariante} fra
le tre pipeline, il che è metodologicamente corretto per confrontarle
fra loro, ma non dice nulla su quanto il modello sia buono nel suo
compito, né quanto costi la quantizzazione a 8 bit.}

\section{Asse 2 --- Costo per pacchetto}

Latenza e throughput. È l'asse su cui il progetto è già forte, ed è
anche quello dove la letteratura è più esigente sul \emph{come} si
misura --- questione rinviata al Capitolo~\ref{ch:come}.

Due precisazioni sulle unità. Primo: la latenza riportata nei lavori
seri è quella del dispositivo sotto test, isolata dal resto, e
accompagnata da un \emph{riferimento} --- tipicamente lo stesso
percorso senza inferenza, per separare il costo del modello dal costo
dell'infrastruttura di forwarding. Secondo: il throughput è una
grandezza \emph{misurata}, non derivata dalla latenza. Sono due cose
diverse, e il Capitolo~\ref{ch:trappole} spiega perché confonderle è
un problema.

I numeri tipici danno la scala del campo: N3IC riporta latenza di
classificazione fino a 100 volte inferiore e throughput da 1,5 a 7
volte superiore rispetto a sistemi software, mantenendo il forwarding
a line rate su NIC a 40\,Gbps~\cite{siracusano22n3ic}.

\section{Asse 3 --- Costo in risorse}

Quanto del dispositivo consuma il modello. La forma concreta dipende
dal bersaglio: su switch programmabili si contano stadi della
pipeline, tabelle match-action, memoria SRAM/TCAM; su FPGA si contano
risorse logiche e blocchi di memoria --- N3IC riporta meno dell'1--2\%
delle risorse di un Virtex7~\cite{siracusano22n3ic}; su eBPF/XDP si
contano istruzioni verificate, dimensione del codice JIT-ato e memoria
delle mappe.

Questo asse ha un ruolo particolare: è l'unico completamente
\emph{deterministico}. Non dipende dal carico, dal rumore o dalla
macchina, il che lo rende l'unica famiglia di metriche confrontabile
fra esperimenti eseguiti in momenti e su hardware diversi.

\section{Asse 4 --- Costo di aggiornamento e flessibilità}

Quanto costa cambiare il modello a sistema in esercizio, e che cosa si
può cambiare senza ricompilare. È l'asse più recente e quello meno
standardizzato, ma è centrale nei lavori che si presentano come
\emph{framework} anziché come singola implementazione:
Planter~\cite{zheng24planter} nasce esattamente per rendere rapido il
ciclo modello-bersaglio-misura, e la sua tesi è che la velocità di
prototipazione sia essa stessa un risultato.

\verdetto{Che cosa comporta per noi}{Questo è l'asse su cui il
progetto ha il contributo più originale. Il confronto fra le tre
pipeline \emph{è} una misura su questo asse, ed è già quantificata:
ricompilazione contro scrittura di una mappa, con circa 926\,ms
misurati per la hardcoded contro frazioni di millisecondo per le altre
due. Va valorizzato, non dato per scontato.}

% ==================================================================
\chapter{Con quali strumenti}
\label{ch:come}

\section{Il protocollo standard: RFC 2544}

Per il costo per pacchetto esiste uno standard, ed è vecchio e
consolidato: la RFC 2544~\cite{rfc2544}, metodologia di benchmarking
per dispositivi di interconnessione di rete. È la base su cui sono
costruiti praticamente tutti gli strumenti commerciali di test.

Definisce poche cose, in modo stretto:

\begin{itemize}
  \item \metrica{Throughput} --- non ``quanti pacchetti al secondo
  riesco a spingere'', ma il \emph{massimo frame rate senza perdita di
  frame}. Si determina per ricerca binaria: si invia a un certo tasso,
  si contano i frame persi, si aggiusta.
  \item \metrica{Latency} --- misurata secondo le definizioni della
  RFC 1242, distinguendo dispositivi store-and-forward da bit-forwarding.
  \item \metrica{Frame loss rate} --- la percentuale di perdita a tassi
  di ingresso crescenti.
  \item \metrica{Back-to-back frames} --- il comportamento in presenza
  di raffiche.
\end{itemize}

Due prescrizioni della RFC meritano attenzione perché toccano
direttamente il nostro caso. La prima: ogni prova dovrebbe durare
\textbf{almeno 120 secondi}. La seconda: la prova va ripetuta
\textbf{almeno 20 volte} e si riporta la media dei valori
registrati~\cite{rfc2544}. Sono requisiti pensati per hardware di
rete, e non vanno trasportati alla cieca su un microbenchmark
kernel-side, ma danno la misura di quanto la comunità consideri
insufficiente una singola prova breve.

Tutte le grandezze vanno inoltre riportate per \emph{diverse dimensioni
di frame} --- tipicamente da 64 byte in su --- perché il costo per
pacchetto e il costo per byte si comportano in modo opposto.

\section{I generatori di traffico}

La conseguenza pratica della RFC 2544 è che serve un generatore capace
di saturare il collegamento e di contare con precisione. Il lavoro di
riferimento su XDP~\cite{hoiland18xdp} --- che è la base tecnologica
di questo progetto --- usa il generatore \textbf{TRex} su schede
Mellanox ConnectX-5 a 100\,Gbps, con CPU Xeon E5-1650 v4. Gli altri
nomi ricorrenti nella letteratura sono MoonGen e pktgen-DPDK.

Il punto non è lo strumento specifico, ma il fatto che la misura del
throughput richiede \emph{traffico vero contro il dispositivo vero}.
Non esiste in letteratura una misura di throughput accettata che sia
stata derivata da un microbenchmark.

\section{Il microbenchmark kernel-side: \texttt{BPF\_PROG\_TEST\_RUN}}

Questo progetto usa \texttt{BPF\_PROG\_TEST\_RUN}, la primitiva che
esegue un programma eBPF già caricato su un pacchetto sintetico
direttamente nel kernel, restituendo il valore di ritorno e una misura
di durata. È una scelta legittima e ben motivata nella tesi: aggira il
problema della consegna end-to-end nell'emulazione.

Va però documentata per quello che è, perché ha limiti precisi. La
documentazione del kernel è esplicita: in questa modalità i programmi
\textbf{non hanno effetti collaterali}; in particolare i pacchetti non
vengono realmente rediretti né scartati, e il valore di ritorno viene
semplicemente riportato allo spazio utente~\cite{kerneldocs}.

Questo significa che sotto \texttt{BPF\_PROG\_TEST\_RUN}:

\begin{itemize}
  \item il costo di \texttt{bpf\_redirect} non viene pagato --- non si
  attraversa driver né coda di trasmissione;
  \item non esiste una NIC, quindi niente DMA, niente interrupt,
  niente pressione sulla cache dovuta al traffico reale;
  \item non c'è perdita di pacchetti, quindi il throughput
  \emph{a perdita nulla} della RFC 2544 non è nemmeno definibile.
\end{itemize}

È un ottimo strumento per confrontare varianti dello stesso programma
in condizioni identiche --- che è esattamente l'uso che ne fa la
Tabella del confronto fra pipeline. Non è uno strumento per affermare
quanto il sistema regga in esercizio.

% ==================================================================
\chapter{Le trappole documentate}
\label{ch:trappole}

Questo capitolo raccoglie tre risultati metodologici pubblicati. Non
sono opinioni: sono lavori che hanno misurato l'errore che si commette
seguendo l'intuizione.

\section{Il bias di misura è la norma, non l'eccezione}

Mytkowicz et al.~\cite{mytkowicz09} hanno mostrato che modificare un
aspetto apparentemente innocuo del setup sperimentale basta a portare
un ricercatore a conclusioni sbagliate. Nei loro esperimenti il bias
di misura si manifesta su tutte le architetture provate, con entrambi
i compilatori provati, e sulla maggior parte dei programmi C di SPEC
CPU2006. Gli effetti misurati sono sufficienti a sovrastimare un
guadagno o addirittura a invertirne il segno.

Il dato più duro del lavoro è bibliometrico: esaminando 133 articoli
recenti da ASPLOS, PACT, PLDI e CGO, nessuno degli articoli con
risultati sperimentali considerava adeguatamente il bias di misura.

Le contromisure che propongono sono due: analisi causale per rilevarlo,
e \emph{randomizzazione del setup} per evitarlo.

\verdetto{Che cosa comporta per noi}{Le nostre latenze sono state
raccolte su una macchina, in una sessione, senza randomizzazione e
senza controllo documentato su frequency scaling, pinning e
isolamento delle CPU. La dispersione osservata --- fino al 178\% fra
minimo e massimo nell'ultima esecuzione --- è coerente con un ambiente
rumoroso.}

\section{Un microbenchmark più veloce non è un sistema più veloce}

Un'analisi recente specificamente su eBPF~\cite{eunomia} sostiene che
quando un'ottimizzazione rende un ciclo stretto di
\texttt{BPF\_PROG\_TEST\_RUN} più veloce del 18\%, questa non è
evidenza sufficiente per abilitarla: l'ottimizzazione può essere del
tutto corretta e restare una cattiva scelta in produzione.

I meccanismi elencati sono concreti: un'ottimizzazione può accelerare
un ciclo aritmetico ma aumentare la dimensione del codice fino a
danneggiare la cache istruzioni in un programma più grande; può
giovare a un backend JIT e non fare nulla su un altro; può ottimizzare
un programma che pesa quasi nulla sul tempo totale dell'applicazione.
E soprattutto: i benchmark isolati misurano singole operazioni, mentre
i programmi reali combinano parsing, lookup su mappe, salti e
riscrittura del pacchetto in un'unica catena, con caratteristiche che
dipendono dall'interazione fra queste operazioni.

\verdetto{Che cosa comporta per noi}{È una descrizione precisa della
nostra situazione. La nostra Pipeline~1 ha 997 istruzioni statiche di
cui ne esegue una frazione, e la tesi lo argomenta già bene. Ma il
corollario vale anche al contrario: il vantaggio della hardcoded
misurato in sandbox potrebbe ridursi sotto traffico reale, dove il
costo del redirect e la pressione sulla cache sono gli stessi per
tutte e tre.}

\section{Ambienti astratti non predicono ambienti reali}

Il terzo risultato viene dal lato routing, ed è il più vicino al
nostro scenario. Boltres et al.~\cite{boltres24} studiano
l'ottimizzazione del routing appresa e trovano che i modelli di rete
``fluidi'' --- quelli astratti, che non modellano l'interazione a
livello di pacchetto --- non sono adeguati per adattamenti su scala di
millisecondi. La conclusione sperimentale è netta: le strategie
apprese addestrate in ambienti fluidi \textbf{non generalizzano} al
setup più realistico a livello di pacchetto.

Osservano inoltre che gli ambienti di valutazione standardizzati nati
per favorire la riproducibilità nella ricerca sul traffic engineering
restano limitati proprio perché usano modelli di rete astratti.

\verdetto{Che cosa comporta per noi}{Il nostro modello decide un
reroute su scala di pacchetto. Valutarlo esclusivamente in un sandbox
che non consegna pacchetti è la versione eBPF dello stesso problema.}

% ==================================================================
\chapter{Dove siamo}
\label{ch:noi}

\section{Quadro sinottico}

La tabella seguente confronta, asse per asse, quello che la letteratura
riporta con quello che il progetto produce oggi.

\begin{longtable}{p{4.2cm}p{4.6cm}p{4.6cm}}
\toprule
\textbf{Metrica} & \textbf{Stato nel progetto} & \textbf{Giudizio} \\
\midrule
\endhead
Accuratezza del modello (float) & assente & \textbf{manca} \\
Accuratezza quantizzata int8 & assente & \textbf{manca} \\
Delta di quantizzazione & assente & \textbf{manca} \\
\midrule
Equivalenza fra pipeline & riferimento intero in Python, confronto su classe scelta & \textbf{solido}, oltre la media \\
Architetture alternative & 65-8-7, 65-4-4-4-7, multi-modello concorrente & \textbf{solido} \\
\midrule
Latenza (microbenchmark) & minimo su 7--15 trial, p50, max, dispersione & \textbf{buono} ma campione esiguo \\
Latenza (end-to-end) & assente & \textbf{manca} \\
Throughput misurato & assente; il valore riportato è derivato da $1/\text{latenza}$ & \textbf{da correggere} \\
Throughput a perdita nulla & assente & \textbf{manca} \\
Sweep sulle dimensioni di frame & assente & \textbf{manca} \\
\midrule
Istruzioni eBPF (xlated) & misurate & \textbf{solido} \\
Codice jited & misurato & \textbf{solido} \\
Letture di mappa per pacchetto & misurate con build strumentata & \textbf{solido} \\
Memoria delle mappe & misurata dal kernel & \textbf{solido} (dopo correzione) \\
\midrule
Costo di aggiornamento modello & misurato con syscall reali & \textbf{solido, originale} \\
Flessibilità architetturale & dimostrata su tre pipeline & \textbf{solido, originale} \\
\midrule
Reroute su guasto & 1 caso su 30 basta a passare & \textbf{debole} \\
Confronto con routing tradizionale & assente & \textbf{manca} \\
Controllo del rumore di misura & non documentato & \textbf{manca} \\
\bottomrule
\end{longtable}

\section{Dove siamo forti}

Va detto con chiarezza, perché è un risultato reale. Il progetto è
\textbf{sopra} la media della letteratura su tre punti.

Il primo è il criterio di correttezza. Confrontare la classe scelta
dal programma eBPF contro un riferimento intero indipendente scritto
in Python, che replica la stessa aritmetica, è più forte di quanto
fanno molti lavori pubblicati, che si limitano a riportare
un'accuratezza aggregata. Qui si verifica l'\emph{equivalenza
esatta}, decisione per decisione.

Il secondo è la trasparenza sulle metriche deterministiche. Distinguere
la dimensione statica del programma dal percorso effettivamente
eseguito, e spiegare perché dividere istruzioni per latenza darebbe un
IPC impossibile, è esattamente il tipo di onestà che il lavoro di
Mytkowicz lamenta essere assente.

Il terzo è l'asse 4. Il confronto costo-per-pacchetto contro
costo-di-aggiornamento, misurato su tre punti di uno spazio di
progetto reale, è un contributo che la letteratura tratta in modo meno
sistematico.

\section{Dove siamo scoperti}

\subsection{Non sappiamo se il modello è buono}

È la lacuna principale. Il progetto dimostra in modo convincente che
le tre pipeline calcolano \emph{la stessa cosa}, ma non dimostra mai
che quella cosa sia \emph{giusta}. Un modello che sceglie sempre la
classe~0 supererebbe tutti i test di equivalenza attuali.

Il test di reroute è l'unico che tocca il merito, ed è debole: passa
se \emph{almeno un caso} su trenta cambia l'uscita. Non verifica che il
cambiamento sia quello corretto, perché non esiste una verità di
riferimento contro cui confrontarlo.

\subsection{Il throughput riportato non è un throughput}

Il valore in Mpps nelle tabelle è calcolato come $1000/\text{latenza in
ns}$. È una \emph{proiezione} della latenza, non una misura: assume
implicitamente serializzazione perfetta, nessuna perdita, nessun costo
di I/O. In particolare, sotto \texttt{BPF\_PROG\_TEST\_RUN} il redirect
non avviene, quindi il costo maggiore del percorso reale è escluso.
Chiamarlo throughput invita un confronto con i numeri della
letteratura che non regge, perché quelli sono ottenuti a perdita nulla
con traffico reale.

\subsection{Due feature del modello sono inerti}

Il progetto documenta già, correttamente, che \texttt{ingress\_iface}
non contribuisce: gli ifindex reali dei nodi Kathara non corrispondono
né alla tabella della Pipeline~1 né all'intervallo accettato da
Pipeline~2 e~3. A questa va aggiunta la feature \texttt{node}, che nel
datapath è indicizzata da \texttt{model\_id} anziché da un
identificatore di nodo: con un solo modello registrato produce sempre
lo stesso contributo su ogni nodo.

Sono 58 delle 65 componenti dell'input vector --- 6 di
\texttt{ingress\_iface} e 52 di \texttt{node} --- che nel laboratorio
non portano informazione. Il modello decide di fatto su
\texttt{link\_state} e TTL. Questo va detto, perché cambia
l'interpretazione di qualunque risultato di accuratezza.

\subsection{Una sola macchina, una sola sessione}

Nessuna randomizzazione del setup, nessun controllo documentato su
governor della frequenza, pinning, isolamento delle CPU o
iperthreading. Con dispersioni fino al 178\% e un campione di 7 trial,
i confronti di latenza fra pipeline vicine non sono difendibili.

% ==================================================================
\chapter{Che cosa fare}
\label{ch:fare}

Il piano è ordinato per rapporto fra valore e costo. I primi due punti
sono quelli che cambiano la tesi; gli altri la rafforzano.

\section{Priorità 1 --- Misurare la qualità del modello}

È la lacuna più grave ed è anche, paradossalmente, la più economica da
colmare: non serve il laboratorio, serve solo Python, e i pezzi
esistono già nel repository.

\begin{enumerate}
  \item \textbf{Costruire un insieme di test con verità di
  riferimento.} Per la topologia Germany50, generare scenari di guasto
  (fino ai 10 guasti simultanei già dichiarati come limite) e calcolare
  per ciascuno il next-hop corretto verso la destinazione con un
  algoritmo classico di cammini minimi. Questa è la verità di
  riferimento che oggi manca.

  \item \textbf{Riportare l'accuratezza su tre livelli:} il modello
  PyTorch in virgola mobile; il riferimento intero in Python; il
  programma eBPF. Il primo confronto misura il costo della
  quantizzazione a 8 bit; il secondo verifica che l'implementazione nel
  kernel non introduca deviazioni. Il secondo confronto il progetto lo
  fa già --- va solo riportato come accuratezza anziché come test
  binario.

  \item \textbf{Dichiarare la baseline.} L'accuratezza va confrontata
  con quella di una politica banale, ad esempio ``scegli sempre il
  next-hop del cammino minimo senza guasti''. Un modello che non batte
  la politica banale non ha valore, e questo confronto è ciò che rende
  il numero interpretabile.
\end{enumerate}

\verdetto{Criterio di accettazione}{Una tabella con accuratezza float,
accuratezza int8, delta di quantizzazione, e accuratezza della
politica banale, su un insieme di test dichiarato e riproducibile.}

\section{Priorità 2 --- Rendere onesto il throughput}

Due opzioni, non alternative.

La via minima, a costo quasi nullo: \textbf{rinominare la metrica}. Non
``Throughput (Mpps)'' ma ``Tasso di inferenza derivato dalla latenza
(Mpps)'', con una nota che dichiara che sotto
\texttt{BPF\_PROG\_TEST\_RUN} il redirect non viene eseguito. Questo
elimina il problema di onestà senza nuove misure.

La via completa, se il tempo lo permette: \textbf{una misura di
throughput vera} su due macchine collegate, o anche su una sola con due
interfacce e namespace separati, con un generatore che saturi il
collegamento. Non serve TRex a 100\,Gbps: bastano \texttt{pktgen} del
kernel o \texttt{trafgen}, purché si riporti il throughput a perdita
nulla e almeno tre dimensioni di frame. Anche a 1\,Gbps il risultato
sarebbe un dato reale, e la differenza fra le tre pipeline sotto carico
è esattamente ciò che la letteratura si aspetta di vedere.

\verdetto{Criterio di accettazione}{O la metrica è rinominata e
annotata, oppure esiste una misura a perdita nulla con generatore
dichiarato, durata dichiarata e almeno tre dimensioni di frame.}

\section{Priorità 3 --- Chiudere il ciclo end-to-end}

La tesi dichiara già che il traffico UDP sulla porta 9999 non viene
consegnato oltre il singolo salto, e che per questo la verifica usa
\texttt{BPF\_PROG\_TEST\_RUN}. È una giustificazione onesta ma resta
una debolezza: il meccanismo di fast-reroute non è mai stato mostrato
funzionare in rete.

Serve almeno \emph{una} dimostrazione end-to-end: un percorso di pochi
salti, con le pipeline attaccate, un flusso da \texttt{h\_src} a
\texttt{h\_dst}, un link fatto cadere a metà flusso, e la prova che il
traffico cambia uscita senza intervento del control plane. Non deve
essere una misura di prestazioni --- basta che sia una prova di
funzionamento. Vale più di dieci tabelle di microbenchmark, perché è
l'unica evidenza diretta della tesi centrale del lavoro.

\section{Priorità 4 --- Irrobustire le misure di latenza}

Interventi a basso costo, tutti documentabili in due righe di
metodologia:

\begin{itemize}
  \item portare i trial da 7 a almeno 30, riportando mediana e
  intervallo di confidenza oltre al minimo;
  \item fissare il governor della frequenza su \texttt{performance} e
  documentarlo;
  \item pinnare il processo su una CPU isolata e dichiarare quale;
  \item randomizzare l'ordine delle pipeline fra le ripetizioni, che è
  la contromisura proposta da Mytkowicz et al.~\cite{mytkowicz09};
  \item ripetere su almeno due sessioni separate da un riavvio.
\end{itemize}

La scelta del minimo come stimatore resta difendibile in ambiente
rumoroso~\cite{chen16}, ma va affiancata da una misura di dispersione
--- cosa che il progetto già fa --- e da un campione adeguato, cosa che
oggi non fa.

\section{Priorità 5 --- Riparare o dichiarare le feature inerti}

Due strade, entrambe accettabili, ma una va scelta esplicitamente.

Ripararle: risolvere gli ifindex reali all'avvio e fornire a
Pipeline~2 e~3 una mappa da ifindex a porta logica; seminare su ogni
nodo il proprio identificatore e usarlo per la feature \texttt{node}
anziché \texttt{model\_id}.

Dichiararle: scrivere nella tesi che nel laboratorio attuale il modello
decide su \texttt{link\_state} e TTL, e che le altre due feature sono
predisposte ma inerti. È una limitazione perfettamente accettabile in
una tesi, purché sia dichiarata e non lasciata al lettore da scoprire.

La scelta peggiore è la terza: riportare l'accuratezza senza
menzionarle.

% ==================================================================
\chapter{In sintesi}

Il progetto misura bene il \emph{costo} del machine learning nel data
plane, su tre punti di uno spazio di progetto, con metriche
deterministiche verificate e un criterio di equivalenza più rigoroso
della media. Questo è il contributo, ed è solido.

Quello che manca è il \emph{valore}: nessuna misura dice se il modello
svolga bene il compito, quanto costi la quantizzazione, né se il
meccanismo funzioni in rete. La letteratura considera queste tre cose
obbligatorie --- non come contorno, ma come prima tabella.

La buona notizia è che la lacuna più grave è anche la più economica da
colmare. Costruire un insieme di test con verità di riferimento e
riportare tre livelli di accuratezza non richiede il laboratorio, non
richiede nuovo hardware, e riusa il riferimento Python già scritto.
È il singolo intervento con il rapporto valore/costo più alto
dell'intero elenco.

% ==================================================================
\begin{thebibliography}{99}

\bibitem{zheng24survey}
C.~Zheng, X.~Hong, D.~Ding, S.~Vargaftik, Y.~Ben-Itzhak, N.~Zilberman,
\emph{In-Network Machine Learning Using Programmable Network Devices:
A Survey}, IEEE Communications Surveys \& Tutorials, vol.~26, no.~2,
2024. DOI: 10.1109/COMST.2023.3344351.

\bibitem{parizotto23survey}
R.~Parizotto, B.~L.~Coelho, D.~C.~Nunes, I.~Haque, A.~Schaeffer-Filho,
\emph{Offloading Machine Learning to Programmable Data Planes: A
Systematic Survey}, ACM Computing Surveys, vol.~56, no.~1, art.~18,
2023. DOI: 10.1145/3605153.

\bibitem{siracusano22n3ic}
G.~Siracusano, S.~Galea, D.~Sanvito, M.~Malekzadeh, G.~Antichi,
P.~Costa, H.~Haddadi, R.~Bifulco,
\emph{Re-architecting Traffic Analysis with Neural Network Interface
Cards}, 19th USENIX Symposium on Networked Systems Design and
Implementation (NSDI~'22), 2022.
\url{https://www.usenix.org/conference/nsdi22/presentation/siracusano}

\bibitem{zheng24iisy}
C.~Zheng, Z.~Xiong, T.~T.~Bui, S.~Kaupmees, R.~Bensoussane,
A.~Bernabeu, S.~Vargaftik, Y.~Ben-Itzhak, N.~Zilberman,
\emph{IIsy: Hybrid In-Network Classification Using Programmable
Switches}, IEEE/ACM Transactions on Networking, vol.~32, no.~3, 2024.
DOI: 10.1109/TNET.2024.3364757.

\bibitem{zheng24planter}
C.~Zheng, M.~Zang, X.~Hong, L.~Perreault, R.~Bensoussane,
S.~Vargaftik, Y.~Ben-Itzhak, N.~Zilberman,
\emph{Planter: Rapid Prototyping of In-Network Machine Learning
Inference}, ACM SIGCOMM Computer Communication Review, vol.~54, no.~1,
pp.~2--21, 2024. DOI: 10.1145/3687230.3687232.

\bibitem{hoiland18xdp}
T.~Høiland-Jørgensen, J.~D.~Brouer, D.~Borkmann, J.~Fastabend,
T.~Herbert, D.~Ahern, D.~Miller,
\emph{The eXpress Data Path: Fast Programmable Packet Processing in
the Operating System Kernel}, 14th International Conference on
emerging Networking EXperiments and Technologies (CoNEXT~'18), 2018.
DOI: 10.1145/3281411.3281443.

\bibitem{rfc2544}
S.~Bradner, J.~McQuaid,
\emph{Benchmarking Methodology for Network Interconnect Devices},
RFC~2544, IETF, 1999. \url{https://www.rfc-editor.org/rfc/rfc2544.html}

\bibitem{mytkowicz09}
T.~Mytkowicz, A.~Diwan, M.~Hauswirth, P.~F.~Sweeney,
\emph{Producing Wrong Data Without Doing Anything Obviously Wrong!},
14th International Conference on Architectural Support for Programming
Languages and Operating Systems (ASPLOS~'09), 2009.
DOI: 10.1145/1508244.1508275.

\bibitem{boltres24}
A.~Boltres, N.~Freymuth, P.~Jahnke, H.~Karl, G.~Neumann,
\emph{Learning Sub-Second Routing Optimization in Computer Networks
requires Packet-Level Dynamics}, Transactions on Machine Learning
Research (TMLR), 2024. arXiv:2410.10377.

\bibitem{chen16}
J.~Chen, J.~Revels,
\emph{Robust benchmarking in noisy environments}, arXiv:1608.04295,
2016.

\bibitem{kerneldocs}
The Linux Kernel documentation,
\emph{Running BPF programs from userspace}.
\url{https://docs.kernel.org/bpf/bpf_prog_run.html}

\bibitem{eunomia}
eunomia,
\emph{When Is an eBPF Optimization Result Strong Enough to Ship?}
\url{https://eunomia.dev/research/ebpf-optimization-evidence-contract/}

\end{thebibliography}

\end{document}
